%% =====================================================================
\section{Introduction}\label{sec:intro}
%% =====================================================================
\IEEEPARstart{M}{any} Internet of Things (IoT) devices no longer send what they sense. A camera trap reports that it saw an animal, a nanosatellite reports that a pixel block contains a ship, and an unmanned aerial vehicle (UAV) that surveys a flood reports that the area below it is under water. Each of these devices runs a small classifier on board and transmits its decision, a label of a few bytes, instead of an image of tens of kilobytes~\cite{shamsoshoara2021flame,kyrkou2020aider,hosny2026tinysem}. For a battery-powered device on a weak wireless link, this saves most of the energy and airtime that a raw upload would cost.

The saving comes at a price. A classifier small enough to run on a UAV or a microcontroller makes mistakes, and it makes them most often on ambiguous inputs: a flooded field under cloud shadow, a dry field after rain. Once the transmitted content is an inference, a delivered update can be wrong. The receiver then holds a fresh but wrong estimate, and it keeps it until a later update corrects it. The age of information (AoI) does not register this error, because it resets on every delivery whether the content is right or wrong. The age of incorrect information (AoII)~\cite{maatouk2020aoii} grows while the receiver's estimate differs from the true state, so it penalizes wrong estimates as well as stale ones. It is the natural measure of how well the receiver knows the scene.

Existing AoII work assumes, however, that the transmitter observes the source through a sensor~\cite{maatouk2020aoii,zakeri2023aoiieh,ngo2025eh}. Such a transmitter knows whether the receiver is wrong and can decide to update exactly when it is. A device whose sensor is a classifier has no such knowledge. It sees only its own label and how confident the classifier is about it. A confident label that disagrees with the receiver probably means that the receiver is wrong. An unsure label could mean either, and sending it may turn a correct estimate into a wrong one. The device has a third option: it can upload the raw data and let a larger classifier on the ground decide, at a much higher energy cost that depends on the link.

Consider a UAV that maps a flood for a ground station (GS) that coordinates rescue teams. If a frame looks clearly flooded and the GS status says dry, the UAV should send its label. If the frame is ambiguous, it should upload the frame when the link is good enough, and wait otherwise. If its label agrees with the GS status, it should stay silent. This paper turns these rules of thumb into a policy with guarantees. Four difficulties stand in the way. First, the UAV never observes whether the GS is wrong. Second, the confidence of a small classifier is itself noisy, and its errors are correlated across consecutive frames because consecutive footprints overlap. Third, the delivery probability and energy of each action depend on the distance to the GS and on blockage. Fourth, the UAV has a fixed energy budget for the whole mission, so a transmission now reduces what is available later.

The UAV keeps a belief over the current AoII of the GS status, updated with the calibrated confidence of each frame and with the acknowledgments of the GS. From the belief we compute in closed form how much each action reduces the expected AoII of the next slot, and a Lyapunov drift-plus-penalty rule weighs this reduction against the energy cost, with a price that rises when the UAV overspends. The main contributions are as follows.
\begin{enumerate}
\item \emph{Model.} We formulate confidence-aware status updating as a constrained partially observable Markov decision process (POMDP). It couples a calibrated per-sample error model of the on-board classifier, a ground classifier, an intermittent air-to-ground link and an energy budget. We show that a belief over the current AoII is a sufficient statistic, and we derive its update, including a prior correction and a tempered likelihood for errors that are correlated in time (Sections~\ref{sec:model} and~\ref{sec:belief}).
\item \emph{Structure.} We derive the one-step expected AoII of each action in closed form. We prove that a label that agrees with the GS has no value and that the per-slot decision has a threshold structure in the confidence: the UAV stays silent at low confidence, uploads the raw frame at intermediate confidence, and sends the label at high confidence (Section~\ref{sec:structure}).
\item \emph{Policy.} We build the online policy CA-DPP on these results and prove that it meets the energy budget on every sample path. CA-DPP needs no knowledge of the confidence distribution, the blockage statistics or the trajectory, and its per-slot cost is linear in the size of the belief (Section~\ref{sec:lyap}).
\item \emph{Evaluation.} We evaluate CA-DPP on inference traces from the Sen1Floods11 dataset~\cite{bonafilia2020sen1floods11} under chip-level cross-validation, against policies adapted from four recent designs~\cite{zakeri2023aoiieh,gan2025taoi,pehlivanoglu2026orbit,basnayaka2025aomi}. We identify when the per-frame confidence pays off: under tight energy budgets and with weak on-board classifiers. We also identify when it does not: with an accurate classifier and ample energy, simpler policies perform as well (Sections~\ref{sec:results} and~\ref{sec:discussion}).
\end{enumerate}

The rest of the paper is organized as follows. Section~\ref{sec:related} reviews related work. Section~\ref{sec:model} presents the system model. Section~\ref{sec:problem} formulates the problem and derives the belief of the UAV about the AoII. Section~\ref{sec:structure} derives the value of each action and the structure of the per-slot decision. Section~\ref{sec:lyap} presents CA-DPP and its energy guarantee. Section~\ref{sec:results} reports the evaluation, Section~\ref{sec:discussion} discusses its implications and limitations, and Section~\ref{sec:concl} concludes. The proofs are in the appendix.
